\chapter{Architettura proposta}
\label{cap:sdd-architettura-proposta}

\section{Panoramica}

Tech4All è organizzato come un'applicazione client/server. Il client è
un'applicazione Next.js che gira nel browser; il server è un processo Node.js
che espone un'interfaccia HTTP con rappresentazioni JSON e conserva i dati in
un DBMS relazionale. Le due parti comunicano solo attraverso quell'interfaccia:
il client non condivide codice con il server e non ne conosce il modello a
oggetti.

Lo stile architetturale è a livelli, replicato identico dentro ogni
sottosistema (\id{DG\_4}):

\begin{description}[leftmargin=1.8cm,style=nextline]
  \item[Rotte] traducono una richiesta HTTP in una chiamata a un servizio.
    Verificano la forma dei dati, applicano il controllo degli accessi e
    serializzano la risposta. Non contengono regole di dominio.
  \item[Servizi] realizzano le regole di dominio, coordinano più DAO e
    delimitano le transazioni. Non conoscono HTTP.
  \item[DAO] traducono fra entità e tabelle. Non contengono decisioni.
  \item[Entità] rappresentano i concetti del dominio e ne custodiscono lo
    stato. Non dipendono da nulla.
\end{description}

La direzione delle dipendenze è sempre la stessa: ogni livello conosce quello
sottostante e ignora quello soprastante.

\section{Decomposizione in sottosistemi}
\label{sec:sdd-sottosistemi}

Il sistema è decomposto in sei sottosistemi applicativi, in corrispondenza
diretta con le aree funzionali del RAD, più due sottosistemi di
infrastruttura attraversati da tutti gli altri.

\diagramma{0.72}{cmp-sottosistemi}{Decomposizione in sottosistemi e loro
dipendenze.}{fig:sdd-cmp}

\begin{longtable}{@{}L{3.4cm} L{11.0cm}@{}}
\toprule
\textbf{Sottosistema} & \textbf{Responsabilità} \\
\midrule
\endfirsthead
\toprule
\textbf{Sottosistema} & \textbf{Responsabilità} \\
\midrule
\endhead
\bottomrule
\endlastfoot

Autenticazione &
Registrazione, riconoscimento delle credenziali, apertura e chiusura della
sessione. Non decide chi può fare cosa: stabilisce soltanto chi è il
chiamante. \\

Gestione account &
Consultazione e modifica dei dati anagrafici, cambio password, cancellazione
dell'account, elenco degli utenti per l'amministratore. \\

Gestione tutorial &
Catalogo, ricerca, filtro, pubblicazione, aggiornamento e ritiro dei
contenuti; sanificazione dell'HTML e normalizzazione delle immagini. \\

Gestione quiz &
Definizione dei quiz, erogazione priva di soluzioni, correzione e
registrazione degli svolgimenti. \\

Gestione feedback &
Inserimento, modifica, rimozione e moderazione delle valutazioni. \\

Gestione obiettivi &
Definizione degli obiettivi e assegnazione automatica dei conseguimenti. \\

\midrule

Middleware di sicurezza &
Verifica del token di sessione, controllo del ruolo, validazione della forma
delle richieste, gestione degli errori e degli upload. \\

Accesso ai dati &
Connessioni al DBMS, esecuzione delle interrogazioni, delimitazione delle
transazioni. \\

\end{longtable}

L'accoppiamento fra sottosistemi applicativi è quasi nullo: l'unica dipendenza
è quella di \emph{Gestione quiz} verso \emph{Gestione obiettivi}, che riflette
la relazione di estensione fra \id{UC\_09} e \id{UC\_10}. Non è una
dipendenza accidentale ma la traduzione di un fatto di dominio: il
conseguimento di un badge è una conseguenza del superamento di un quiz.

\section{Mappatura hardware/software}
\label{sec:sdd-deployment}

Il sistema è distribuito su tre nodi logici: il dispositivo dell'utente, il
nodo applicativo e il nodo dati. In sviluppo i due nodi server coincidono; la
separazione è logica e non impone una topologia.

\diagramma{0.56}{dep-deployment}{Mappatura hardware/software.}{fig:sdd-dep}

Il server non conserva stato in memoria fra una richiesta e l'altra: la
sessione vive interamente nel token firmato che il browser riallega a ogni
richiesta. È la proprietà che rende il nodo applicativo replicabile senza
alcuna condivisione di stato, se in futuro servisse.

I file caricati risiedono sul filesystem del nodo applicativo e sono serviti
staticamente. È la scelta più semplice compatibile con \id{DG\_6}; il prezzo
è che, in presenza di più nodi applicativi, quella cartella dovrebbe essere
condivisa.

\section{Gestione dei dati persistenti}
\label{sec:sdd-persistenza}

\subsection{Scelta della tecnologia}

La persistenza è affidata a MySQL 8. La scelta discende dalla natura dei
dati, non da preferenze: il dominio è fatto di entità con relazioni fisse e
molteplicità note, e gli invarianti più importanti — l'unicità di un feedback
per coppia utente-tutorial, l'esistenza di un solo quiz per tutorial, la
rimozione a cascata di ciò che dipende da un tutorial — sono esprimibili come
vincoli del modello relazionale. Affidarli al DBMS significa che restano veri
anche se il codice applicativo sbaglia (\id{DG\_5}).

\subsection{Schema logico}

\diagrammaLargo{er-database}{Schema logico della base di dati.}{fig:sdd-er}

Le tre tabelle \file{feedback}, \file{svolgimento} e \file{conseguimento}
hanno chiave primaria composta dalle due chiavi esterne che collegano: è la
traduzione diretta delle classi di associazione del modello di analisi, e
rende impossibile per costruzione l'esistenza di un duplicato.

\subsection{Vincoli affidati al database}

\begin{longtable}{@{}L{5.4cm} L{9.0cm}@{}}
\toprule
\textbf{Vincolo} & \textbf{Realizzazione} \\
\midrule
\endfirsthead
\toprule
\textbf{Vincolo} & \textbf{Realizzazione} \\
\midrule
\endhead
\bottomrule
\endlastfoot

Un utente lascia al più un feedback per tutorial &
Chiave primaria \id{(utente\_id, tutorial\_id)} su \file{feedback} \\

Un tutorial ha al più un quiz &
Vincolo di unicità su \file{quiz.tutorial\_id} \\

Un utente ha al più uno svolgimento per quiz &
Chiave primaria \id{(utente\_id, quiz\_id)} su \file{svolgimento} \\

Rimuovere un tutorial rimuove quiz, domande, risposte, svolgimenti e feedback &
Vincoli di integrità referenziale con \id{ON DELETE CASCADE} \\

La valutazione è compresa fra 1 e 5 &
Vincolo \id{CHECK} sulle tabelle \file{feedback} e \file{tutorial} \\

Il ruolo assume solo i valori previsti &
Tipo \id{ENUM} su \file{utente.ruolo} \\

La valutazione media riflette i feedback ricevuti &
Tre trigger su \file{feedback} che ricalcolano \file{tutorial.valutazione} \\

\end{longtable}

\subsection{Transazioni}

Due operazioni coinvolgono più tabelle e devono essere atomiche:

\begin{itemize}
  \item la \textbf{definizione di un quiz}, che inserisce il quiz, le sue
        domande e le risposte di ciascuna domanda;
  \item la \textbf{correzione di un quiz}, che registra lo svolgimento,
        ricalcola il numero di quiz superati e assegna gli eventuali
        conseguimenti.
\end{itemize}

Entrambe sono racchiuse nella funzione \id{withTransaction}, che apre la
transazione, esegue il lavoro, effettua il commit e in caso di eccezione il
rollback, rilasciando sempre la connessione. I DAO coinvolti accettano la
connessione della transazione come parametro: usare il pool comune
all'interno di una transazione già aperta porterebbe la stessa richiesta ad
attendere i lock che essa stessa detiene.

\diagramma{0.86}{sd-correzione-quiz}{Correzione di un quiz: propagazione della
connessione transazionale ai DAO coinvolti.}{fig:sdd-sd-correzione}

\section{Controllo degli accessi e sicurezza}
\label{sec:sdd-accessi}

\subsection{Riconoscimento del chiamante}

La sessione è un token JWT~\cite{rfc7519} firmato dal server, contenente
l'identificativo dell'utente e il suo ruolo, depositato in un cookie
\id{httpOnly}, \id{SameSite=Lax} e, in produzione, \id{Secure}. Il token non
è leggibile dal JavaScript di pagina: una vulnerabilità di scripting non
permette di esfiltrarlo.

\diagramma{0.66}{sd-autenticazione}{Autenticazione e successiva richiesta
protetta.}{fig:sdd-sd-auth}

\subsection{Matrice di controllo degli accessi}

Il controllo è realizzato da due middleware: \id{requireAuth} rifiuta le
richieste prive di una sessione valida; \id{requireAdmin} rifiuta in più
quelle di chi non ha ruolo amministratore. Sono attraversati prima che il
sottosistema applicativo entri in gioco.

\begin{longtable}{@{}L{6.2cm} C{2.4cm} C{2.4cm} C{2.4cm}@{}}
\toprule
\textbf{Operazione} & \textbf{Visitatore} & \textbf{Utente} & \textbf{Amm.} \\
\midrule
\endfirsthead
\toprule
\textbf{Operazione} & \textbf{Visitatore} & \textbf{Utente} & \textbf{Amm.} \\
\midrule
\endhead
\bottomrule
\endlastfoot

Registrarsi, autenticarsi              & \checkmark & --- & --- \\
Consultare catalogo, tutorial, feedback, quiz, obiettivi & \checkmark & \checkmark & \checkmark \\
Consultare e modificare il proprio profilo & --- & \checkmark & \checkmark \\
Eliminare il proprio account           & --- & \checkmark & \checkmark \\
Consegnare un quiz                     & --- & \checkmark & \checkmark \\
Gestire il proprio feedback            & --- & \checkmark & \checkmark \\
Consultare i propri badge              & --- & \checkmark & \checkmark \\
Pubblicare, aggiornare, ritirare tutorial & --- & --- & \checkmark \\
Definire e rimuovere quiz              & --- & --- & \checkmark \\
Caricare ed eliminare immagini         & --- & --- & \checkmark \\
Definire, aggiornare, eliminare obiettivi & --- & --- & \checkmark \\
Consultare ed eliminare account altrui & --- & --- & \checkmark \\
Moderare i feedback altrui             & --- & --- & \checkmark \\

\end{longtable}

\subsection{Provenienza dell'identità}

Un principio attraversa tutto il sistema (\id{DG\_1}): l'autore di
un'operazione è sempre l'utente della sessione. Nessuna rotta accetta un
identificativo utente dal corpo o dai parametri della richiesta per stabilire
per conto di chi agire. L'unico caso in cui un identificativo altrui compare
in una rotta è la moderazione dei feedback, dove indica il \emph{bersaglio}
dell'operazione e non il suo autore, ed è accessibile solo agli
amministratori.

\subsection{Altre misure}

\begin{longtable}{@{}L{4.2cm} L{10.2cm}@{}}
\toprule
\textbf{Minaccia} & \textbf{Contromisura} \\
\midrule
\endfirsthead
\toprule
\textbf{Minaccia} & \textbf{Contromisura} \\
\midrule
\endhead
\bottomrule
\endlastfoot

Furto di credenziali dal database &
Le password sono conservate come hash bcrypt~\cite{provos1999}, con un fattore
di costo configurabile. \\

Enumerazione degli account &
Credenziali errate ed email inesistente producono la stessa risposta. \\

Iniezione di codice nei contenuti &
L'HTML dei tutorial è sanificato in ingresso contro una whitelist di
marcatori e attributi; schemi di URL diversi da \id{http}, \id{https} e
\id{mailto} sono rimossi~\cite{owasp2021}. \\

Iniezione SQL &
Tutte le interrogazioni usano parametri preparati; nessun valore proveniente
dall'esterno è concatenato nel testo di una query. \\

Traversamento del filesystem &
I nomi dei file caricati sono generati dal server; l'eliminazione accetta solo
nomi semplici, risolti e verificati dentro la cartella consentita. \\

Richieste da origini non previste &
Il CORS autorizza una sola origine, configurata, e ammette le credenziali solo
per essa. \\

Divulgazione di dettagli interni &
Gli errori imprevisti producono una risposta generica; i dettagli restano nei
log del server. \\

\end{longtable}

\section{Controllo del flusso globale}
\label{sec:sdd-flusso}

Il sistema è guidato dagli eventi: non esiste un flusso di controllo proprio,
ma una sequenza di reazioni a richieste HTTP. Ogni richiesta attraversa
sempre le stesse fasi, nello stesso ordine.

\begin{enumerate}
  \item \textbf{Ingresso}: parsing del corpo e dei cookie, verifica
        dell'origine.
  \item \textbf{Autenticazione}: lettura e verifica del token, se la rotta la
        richiede.
  \item \textbf{Autorizzazione}: verifica del ruolo, se la rotta la richiede.
  \item \textbf{Validazione}: controllo della forma dei dati; una violazione
        interrompe qui, senza raggiungere il servizio.
  \item \textbf{Dominio}: il servizio applica le regole e coordina i DAO.
  \item \textbf{Serializzazione}: l'entità è tradotta in DTO.
  \item \textbf{Uscita}: la risposta è emessa; se una fase ha sollevato un
        errore, il gestore centralizzato lo traduce in uno stato HTTP.
\end{enumerate}

Le fasi 2, 3, 4 e 7 sono realizzate da middleware e sono quindi identiche per
tutti i sottosistemi (\id{DG\_2}, \id{DG\_4}). Un solo processo serve tutte
le richieste; la concorrenza è quella del ciclo di eventi di Node.js, e non
esiste stato condiviso in memoria fra richieste diverse.

\section{Condizioni limite}
\label{sec:sdd-condizioni-limite}

\subsection{Avvio}

All'avvio il server verifica la presenza di tutte le variabili d'ambiente
obbligatorie e fallisce immediatamente elencando quelle mancanti, prima di
mettersi in ascolto. È una scelta deliberata: un sistema che parte senza il
segreto di firma o senza le credenziali del database fallirebbe più tardi, in
modo meno comprensibile.

Il pool di connessioni è creato al primo utilizzo e non apre connessioni
all'avvio: l'indisponibilità del database non impedisce al processo di
partire, ma si manifesta come errore sulle richieste che ne hanno bisogno.

\subsection{Spegnimento}

Alla ricezione di \id{SIGINT} o \id{SIGTERM} il server smette di accettare
nuove connessioni, attende il completamento di quelle in corso, chiude il
pool e termina. Le transazioni aperte al momento dell'interruzione vengono
annullate dal DBMS alla caduta della connessione: il database non resta in
uno stato intermedio.

\subsection{Fallimenti}

\begin{longtable}{@{}L{4.6cm} L{9.8cm}@{}}
\toprule
\textbf{Condizione} & \textbf{Comportamento del sistema} \\
\midrule
\endfirsthead
\toprule
\textbf{Condizione} & \textbf{Comportamento del sistema} \\
\midrule
\endhead
\bottomrule
\endlastfoot

Database irraggiungibile &
Le richieste che richiedono dati falliscono con \id{500} e un messaggio
generico; le rotte statiche continuano a rispondere. \\

Errore durante una transazione &
Rollback automatico: nessuna scrittura parziale resta applicata. \\

File di upload non scrivibile &
La pubblicazione fallisce con un errore di validazione; il tutorial non viene
creato. \\

Token scaduto o manomesso &
Trattato come sessione assente: \id{401} sulle rotte protette, accesso
anonimo su quelle pubbliche. \\

Richiesta con corpo eccessivo &
Rifiutata dal limite di dimensione prima di raggiungere l'applicazione. \\

Eccezione non prevista &
Intercettata dal gestore centralizzato, registrata nei log e tradotta in
\id{500} privo di dettagli interni. \\

\end{longtable}
